\chapter{Conclusion}\label{ch:conclusion}

This thesis developed the \textit{Mathematica} package \textit{AntCalc} in an effort 
to make the computation of QCD antenna functions faster than the common from-scratch calculation, 
easier, standardised, and streamlined.

The resulting framework was implemented completely within \textit{Mathematica}, 
using the \textit{FeynCalc} environment (including \textit{FeynArts} and 
\textit{FeynHelpers}) to handle the building stage and
\textit{LiteRed2} for the IBP reduction in the integration stage. That framework was 
shown to produce all massless quark-antiquark final-final antennae through NNLO. 
It was also shown to be able to be extended to the massive regime by implementing the
massive $A_3^0$ route in its entirety. It begins by 
using \textit{FeynArts} to produce one or more Feynman diagrams relevant for the 
process under study and generates their amplitudes. 
It then develops the amplitude into one antenna or multiple 
antenna components using \textit{FeynCalc}. 
The unintegrated antenna is then handed to \textit{LiteRed2} 
for IBP reduction and, finally, after the pre-computed MI are identified and substituted,
the integrated antenna is Laurent-expanded around the dimensional regularisation 
parameter $\epsilon$.

These unintegrated and integrated antennae were shown to be in agreement with 
established results in the literature. Furthermore, the massless $A_3^0$ antenna, 
whose workflow was shown in great detail, was locally validated by computing the 
eikonal factor and the Altarelli--Parisi splitting function, using a method that 
may be used to validate other appropriate antennae. All massless integrated antenna results 
were also validated when we computed the massless hadronic R-ratio. 

The package is beginning to reach the computational limits of its current 
workflow, running completely within Mathematica. As such, the most pressing piece of 
future work is the overhaul to a new, truly modular and multi-language, workflow,
that is able to harness the best characteristics of different tools, like Kira for 
IBP reduction and FORM for algebraic manipulation, without compromising efficiency.
Once this overhaul is done, \textit{AntCalc} could be expanded 
to handle much more complex cases, particularly those in the massive regime, 
other antenna types, such as quark-gluon, other kinematic configurations, such as 
initial-final antennae, and those in higher perturbative orders, such as N$^3$LO.